\subsection{Игры}
\textbf{Как решать.} (1) Свести игру к \emph{независимым} подыграм (куча, ряд, компонента): ходы в одной не влияют на другую. (2) Для каждой подигры посчитать $g$ динамикой (перебор ходов, $\mathrm{mex}$ по значениям Гранди результатов). (3) Взять xor $g$ всех подыгр: ненулевой -- выигрывает ходящий. Если игра на графе с циклами, решаем ретроанализом (позиция проигрышна, если все ходы ведут в выигрышные).
\par \textbf{Шпрага--Гранди:} $g(v)=\mathrm{mex}\{g(u):v\to u\}$; сумма игр -- xor значений; позиция проигрышна $\iff g=0$.
\par \textbf{Ним:} ходит первый и выигрывает $\iff\bigoplus a_i\ne0$. Мизер-ним (проигрывает взявший последним): как ним, если есть куча $>1$; иначе первый выигрывает при \emph{чётном} числе единиц. Лестничный ним: считаются только нечётные ступени. Бэш (берём $1..k$): $g=n\bmod(k{+}1)$. Уайтхоф (две кучи, ходы: из одной или поровну из обеих): P-позиции $(\lfloor k\varphi\rfloor,\lfloor k\varphi^2\rfloor)$, $\varphi=\frac{1+\sqrt5}2$. Кайлс: $g(0..12)=0,1,2,3,1,4,3,2,1,4,2,6,4$ (далее период 12 начиная с $n=71$).
\par Игры на графах без циклов -- динамика по порядку; с циклами -- ретроанализ от терминальных позиций (счётчики выходов). Минимакс: альфа-бета, порядок ходов по оценке.
\subsection{Приёмы и грабли}
\par \textbf{Mo:} блок $n/\sqrt q$, $O(n\sqrt q)$; с обновлениями блок $n^{2/3}$, сортировка по $(\lfloor t/K\rfloor,\lfloor l/K\rfloor,r)$, $O(n^{5/3})$.
\par \textbf{Оптимизации ДП.} Неравенство четырёхугольника для $w$: $w(a,c)+w(b,d)\le w(a,d)+w(b,c)$ ($a\le b\le c\le d$) $\Rightarrow$ монотонность опта: D\&C ($O(n\log n)$ на слой), Кнут для интервалов $opt[l][r{-}1]\le opt[l][r]\le opt[l{+}1][r]$ ($O(n^2)$; нужна ещё монотонность $w$ по вложению). Выпуклость по числу групп $\Rightarrow$ лямбда-оптимизация (бинпоиск штрафа).
\par \textbf{Код Грея:} $g(i)=i\oplus(i\gg1)$, соседние отличаются в бите. \textbf{Битовые:} \texttt{x\&-x} -- младший бит; \texttt{\_\_builtin\_ctz/clz/popcount}; перебор подмасок \texttt{for(s=m;s;s=(s-1)\&m)}; битсет ускоряет в 64 раза.


